\chapter{Motivating Use Case}
\label{sec:use-case}

A current use case in the field of fusion energy science is presented
in~\cite{10678736}, where the objective is to operate 
reactors such as the International Thermonuclear Experimental Reactor
(ITER). These reactors are equipped with sensors to monitor the state of
the plasma and of the reactor itself. The sensor data 
can be processed for several real-time purposes, such as the
steering of experiments between shots of the reactors, or
the detection of catastrophic event precursors.
This data is processed by applications structured as chains
of tasks that perform data compression/decompression, data
reduction/filtering, simulation, analysis, or visualization.  These
applications are executed periodically throughout the lifespan of the
operation of the reactor, with results needed by the end of each
period.  Fig.~\ref{fig:use_case} depicts a three-task chain for the
specific application configuration described in~\cite{9912722}.


\begin{figure}
\begin{center}
\scalebox{0.8}{
  \begin{tikzpicture}[
    node distance=2cm,
    box/.style={rectangle, draw, thick, minimum width=2cm, minimum height=1cm, align=center},
    arrow/.style={-{Stealth[length=2mm]}, thick}
]

% Nodes
	\node[box, align=center] (task1) {$T_1$: Simulation\\with XGC~\cite{10.1063/1.5020792}};
	\node[box, right=of task1, align=center] (task2) {$T_2$: Data Refactoring\\with ADIOS~\cite{GODOY2020100561}};
	\node[box, below=0.7cm of task2, align=center] (task3) {$T_3$: Analysis/Visualization\\with VTK-m~\cite{7466740}};

% Stick figure user
\node[left=2cm of task3, xshift=0.3cm] (user-anchor) {};
\node[left=2cm of task3, yshift=0.4cm] (user) {};
\draw[thick] (user) circle (0.20cm); % head
\draw[thick] (user) ++(0,-0.25) -- ++(0,-0.30); % body
\draw[thick] (user) ++(0,-0.55) -- ++(-0.3,-0.35); % left leg
\draw[thick] (user) ++(0,-0.55) -- ++(0.3,-0.35); % right leg
\draw[thick] (user) ++(0,-0.35) -- ++(-0.2,-0.25); % left arm
\draw[thick] (user) ++(0,-0.35) -- ++(0.2,-0.25); % right arm

% Incoming arrow to Task1
\coordinate[left=1.5cm of task1] (input);
\draw[arrow] (input) -- node[above, align=center] {input\\data} (task1);

% Arrows with labels
\draw[arrow] (task1) -- node[above, align=center] {simulation\\results} (task2);
\draw[arrow] ($(task2)+(-1cm,-0.5cm)$) -- node[right, align=center, yshift=+0.1cm] {reduced results} ($(task3)+(-1cm,+0.5cm)$);
\draw[arrow] (task3) -- node[above, align=center] {Poincar\'e\\maps} (user-anchor);

\end{tikzpicture}

}
\end{center}
\caption{Example three-task chain for a fusion energy science application~\cite{9912722}.}
\label{fig:use_case}
\end{figure}

For fault-tolerance purposes, the authors in~\cite{10678736} propose using
approximate computing techniques to trade off output quality for
shorter execution times. While many such techniques are
available~\cite{10.1145/3716845, 10.1145/3711683}, they focus on data size
reduction using generic compression methods or methods specifically
designed for scientific datasets~\cite{10.1145/3733104}.  Regardless, some
or all of the tasks in the application come with different execution modes
that correspond to, e.g., different compression factors, different data
refactoring options, and different precision levels.  The above use case
provides the motivation for the formulation of the scheduling problem in
the next chapter.






% =========================================================
% 
% See highlights and speech bubbles in the PDFs
% 
% - Error-Controlled Lossy Compression...: BigData18-SZ2.0.pdf
% 	- A paper on lossy stuff for science (cited by paper below)
% 	- May have some data to get some actual values...
% 
% - RAPIDS: ...  (3588195.3592983.pdf):
% 	TLDR:A good ref and source of refs about approximate computing in science.
%              It's basically trying to pick good option for data resilience
% 
% 	- Paragraph/refs about why lossless is not sufficiently scalable
% 	- Paragraph/ref about lossy stiuff and why its good for science
% 	- Paragraph/ref about data refactoring: which is what they use
% 	- Fault-tolerance:
% 		- how to pick data encoding (erasure + multi-level
% 		  refactoring) to maximize fidelity and minimize overhead
% 		  given the probability of losing a storage system
% 
% - Hybrid Analysis...: Hybrid_Analysis.pdf
% 	- Not sure it's relevant at all
% 
% - Towards Resilience Near Real-Time...: Suter_Towards.pdf
%   - THE reference for "timely workflow execution". Two technniques:
% 	- Replicas: nice but costly - Data reduction: faster i/O, faster
% 	  computation 
%         - PROBLEM: Mention a deadline in seconds! And there are faults?  
%         - Mentions that algorithms could also be tweaked but say
% 	  it's better to do data-reduction 
%         - techniques:
% 		- lossy or lossless compression - data reduction (see papers above) 
%                 - perhaps decompress/reconstruct 
% 	- Their problem:
% 		- Deadline 
%                 - Maximum number of failures  (They do R replicas to tolerate R failures) 
%                 - Maximum resource budget
% 	- We can say we consider a more general problem?
% 
